\documentclass[a4paper,12pt]{article}

\usepackage[T2A]{fontenc}
\usepackage[utf8]{inputenc}
\usepackage[russian]{babel}

\usepackage{amsmath,amssymb,mathtools}
\usepackage{helvet}
\renewcommand{\familydefault}{\sfdefault}
\usepackage{icomma}

\usepackage[a4paper,margin=2.2cm,headheight=15pt]{geometry}
\usepackage{microtype}
\usepackage{tikz}
\usepackage{tabularx,booktabs}
\usepackage{enumitem}
\usepackage{hyperref}
\usepackage{parskip}
\usepackage{fancyhdr}
\usepackage{setspace}
\usepackage{xcolor}

\definecolor{accent}{HTML}{008080}
\definecolor{darkaccent}{HTML}{004D4D}
\definecolor{textgray}{HTML}{333333}
\definecolor{softgray}{HTML}{EEF3F3}
\definecolor{muted}{HTML}{6B7777}

\hypersetup{
    colorlinks=true,
    linkcolor=darkaccent,
    urlcolor=darkaccent,
    pdftitle={Чек-лист: производительность и растворы},
    pdfauthor={Лёвин Артём Александрович}
}

\onehalfspacing
\setlength{\parindent}{0pt}
\setlist[itemize]{leftmargin=1.6em,itemsep=0.25em,topsep=0.35em}
\setlist[enumerate]{leftmargin=1.8em,itemsep=0.6em,topsep=0.4em}

\pagestyle{fancy}
\fancyhf{}
\fancyhead[L]{\small\color{muted}Текстовые задачи}
\fancyhead[R]{\small\color{muted}Ксения Клыкова}
\fancyfoot[C]{\small\color{muted}\thepage}
\renewcommand{\headrulewidth}{0.3pt}
\renewcommand{\headrule}{%
  \hbox to\headwidth{\color{accent!45}\leaders\hrule height \headrulewidth\hfill}
}

\newcommand{\topic}[1]{%
    \vspace{0.5em}
    {\Large\bfseries\color{darkaccent}#1}\par
    \vspace{0.25em}
    {\color{accent!55}\rule{\linewidth}{0.6pt}}\par
    \vspace{0.45em}
}

\newcommand{\subtopic}[1]{%
    \vspace{0.4em}
    {\large\bfseries\color{darkaccent}#1}\par
    \vspace{0.15em}
}

\newcommand{\keyidea}[1]{%
    \textbf{\color{darkaccent}Ключевая идея.} #1
}

\begin{document}

% ============================================================
% ТИТУЛЬНЫЙ ЛИСТ
% ============================================================

\begin{titlepage}
\thispagestyle{empty}

\begin{center}

\vspace*{2.6cm}

{\fontsize{18}{22}\selectfont\color{muted}ЧЕК-ЛИСТ\par}

\vspace{0.8cm}

{\fontsize{28}{34}\selectfont\bfseries\color{darkaccent}
Текстовые задачи\par}

\vspace{0.3cm}

{\fontsize{20}{26}\selectfont\bfseries\color{accent}
Производительность и растворы\par}

\vspace{1.4cm}

{\large
ЕГЭ по профильной математике\par}

\vspace{2.0cm}

{\large\bfseries
Лёвин Артём Александрович\par}

\vspace{0.25cm}

{\large
эксклюзивно для Ксении Клыковой\par}

\vfill

{\large \today\par}

\vspace{1.4cm}

\begin{tikzpicture}[x=1cm,y=1cm]
    \draw[
        accent!55,
        line width=1.2pt,
        line cap=round
    ]
    (0,0)
    .. controls (2.0,0.55) and (4.0,0.55) .. (5.5,0)
    .. controls (7.0,-0.55) and (9.0,-0.55) .. (11,0);
\end{tikzpicture}

\vspace{0.5cm}

\end{center}
\end{titlepage}

% ============================================================
% ОСНОВНОЙ ЧЕК-ЛИСТ
% ============================================================

\topic{1. Общая стратегия текстовой задачи}

Практически вся работа строится по одной схеме:

\begin{enumerate}
    \item Определите величины, которые участвуют в задаче.
    \item Составьте таблицу.
    \item Разместите неизвестную величину там, где её проще всего описать.
    \item Заполните известные данные.
    \item Формульный столбец заполните через основную формулу темы.
    \item Переведите единицы измерения к единому виду.
    \item Составьте уравнение по связи между строками таблицы.
    \item Решите уравнение и проверьте смысл найденного значения.
\end{enumerate}

\keyidea{Таблица помогает перевести текст условия на язык формул и увидеть связь между величинами.}

\subtopic{Перед составлением уравнения}

Проверьте:

\begin{itemize}
    \item одинаковы ли единицы времени;
    \item какая величина больше и какая меньше;
    \item какие величины можно складывать;
    \item имеет ли неизвестная физический смысл;
    \item допустимы ли полученные знаменатели.
\end{itemize}

% ============================================================

\topic{2. Производительность, время и результат}

Основные величины:

\begin{center}
\renewcommand{\arraystretch}{1.35}
\begin{tabularx}{0.88\linewidth}{@{} l X c @{}}
\toprule
\textbf{Величина} & \textbf{Смысл} & \textbf{Обозначение} \\
\midrule
Результат работы & выполненная работа, заказ, число деталей, объём воды и т.\,п. & $A$ \\
Производительность & результат за единицу времени & $p$ \\
Время & продолжительность работы & $t$ \\
\bottomrule
\end{tabularx}
\end{center}

Основная формула:

\[
\boxed{A=pt}
\]

Из неё:

\[
\boxed{p=\frac{A}{t}},
\qquad
\boxed{t=\frac{A}{p}}.
\]

Если весь заказ принимается за единицу:

\[
A=1,
\qquad
p=\frac1t.
\]

\subtopic{Совместная работа}

Если несколько исполнителей работают одновременно и их результаты складываются, складываются их производительности:

\[
\boxed{p_{\text{общ}}=p_1+p_2+\ldots}
\]

Для двух исполнителей, выполняющих один заказ:

\[
\boxed{
\frac1{t_{\text{совм}}}
=
\frac1{t_1}+\frac1{t_2}
}.
\]

\keyidea{При совместной работе складываются скорости выполнения работы.}

% ============================================================

\topic{3. Пример: два мастера}

Первый мастер выполняет заказ за $3$ часа, второй --- за $6$ часов.

Пусть совместно они выполнят заказ за $x$ часов.

\begin{center}
\renewcommand{\arraystretch}{1.35}
\begin{tabularx}{0.82\linewidth}{@{} X c c c @{}}
\toprule
& \textbf{Результат} & \textbf{Время} & \textbf{Производительность} \\
\midrule
Первый мастер & $1$ & $3$ & $\dfrac13$ \\
Второй мастер & $1$ & $6$ & $\dfrac16$ \\
Вместе & $1$ & $x$ & $\dfrac1x$ \\
\bottomrule
\end{tabularx}
\end{center}

Совместная производительность:

\[
\frac1x=\frac13+\frac16=\frac12.
\]

Следовательно,

\[
\boxed{x=2\text{ ч}}.
\]

Ограничение:

\[
x>0.
\]

% ============================================================

\topic{4. Если дана производительность}

Петя отвечает на $6$ вопросов в час, Ваня --- на $7$ вопросов в час.
Они выполняют одинаковые тесты, причём Петя заканчивает на $20$ минут позже.

Пусть в тесте $x$ вопросов.

\begin{center}
\renewcommand{\arraystretch}{1.35}
\begin{tabularx}{0.78\linewidth}{@{} X c c c @{}}
\toprule
& \textbf{Производительность} & \textbf{Время} & \textbf{Результат} \\
\midrule
Петя & $6$ & $\dfrac{x}{6}$ & $x$ \\
Ваня & $7$ & $\dfrac{x}{7}$ & $x$ \\
\bottomrule
\end{tabularx}
\end{center}

Поскольку производительность дана в вопросах за час, переводим минуты в часы:

\[
20\text{ мин}=\frac{20}{60}\text{ ч}=\frac13\text{ ч}.
\]

Петя работает дольше:

\[
\frac{x}{6}-\frac{x}{7}=\frac13.
\]

\[
\frac{x}{42}=\frac13,
\qquad
\boxed{x=14}.
\]

\keyidea{Фраза «за час выполняет 6 заданий» задаёт производительность, поскольку указан результат за единицу времени.}

% ============================================================

\topic{5. Сравнение времён и производительностей}

Если два объекта работают одинаковое время, больший результат соответствует большей производительности.

Например, первая труба за одинаковое время наполнила $770$ л, вторая --- $830$ л.

Пусть производительность первой трубы равна $x$ л/мин, а вторая пропускает на $6$ л/мин больше:

\[
p_1=x,
\qquad
p_2=x+6.
\]

Времена:

\[
t_1=\frac{770}{x},
\qquad
t_2=\frac{830}{x+6}.
\]

Поскольку времена равны:

\[
\boxed{
\frac{770}{x}=\frac{830}{x+6}
}.
\]

Ограничения:

\[
x>0,
\qquad
x+6>0.
\]

\subtopic{Полезный ориентир}

При одинаковом времени:

\[
A_2>A_1
\quad\Longrightarrow\quad
p_2>p_1.
\]

Это позволяет правильно определить, где записать $x+d$, если в условии сказано лишь, что одна производительность на $d$ больше другой.

% ============================================================

\topic{6. Разность времён}

Если один исполнитель заканчивает позже другого на $\Delta t$, уравнение обычно имеет вид

\[
\boxed{t_{\text{большее}}-t_{\text{меньшее}}=\Delta t}.
\]

Например:

\[
\frac{x}{6}-\frac{x}{7}=\frac13.
\]

При составлении подобных уравнений сначала определите, кто работает дольше.

\subtopic{Единицы измерения}

Если производительность дана за час, время также выражайте в часах:

\[
15\text{ мин}=\frac14\text{ ч},
\qquad
20\text{ мин}=\frac13\text{ ч},
\qquad
30\text{ мин}=\frac12\text{ ч}.
\]

Если все данные представлены в минутах, можно вести решение целиком в минутах.

% ============================================================

\topic{7. План и фактическая производительность}

Частая модель:

\[
p_{\text{план}}=x,
\qquad
p_{\text{факт}}=x+d.
\]

Если объём работы одинаков:

\[
t_{\text{план}}=\frac{A}{x},
\qquad
t_{\text{факт}}=\frac{A}{x+d}.
\]

Если работа закончена на $\Delta t$ раньше:

\[
\boxed{
\frac{A}{x}-\frac{A}{x+d}=\Delta t
}.
\]

Здесь обязательно:

\[
x>0,
\qquad
x+d>0.
\]

После решения квадратного уравнения выбирается положительное значение, соответствующее смыслу производительности.

% ============================================================

\topic{8. Растворы и смеси}

Основные величины:

\begin{center}
\renewcommand{\arraystretch}{1.35}
\begin{tabularx}{0.9\linewidth}{@{} l X c @{}}
\toprule
\textbf{Величина} & \textbf{Смысл} & \textbf{Обозначение} \\
\midrule
Количество раствора & масса или объём раствора в рамках модели задачи & $m$ \\
Концентрация & доля растворённого вещества & $c$ \\
Количество вещества & количество растворённого вещества & $m_{\text{в-ва}}$ \\
\bottomrule
\end{tabularx}
\end{center}

Главная формула:

\[
\boxed{
m_{\text{в-ва}}=m\,c
}.
\]

Концентрация в вычислениях записывается долей:

\[
15\%=0{,}15,
\qquad
23\%=0{,}23,
\qquad
40\%=0{,}40.
\]

Для перевода результата обратно в проценты:

\[
c=0{,}18
\quad\Longrightarrow\quad
18\%.
\]

\subtopic{Баланс вещества}

При смешивании:

\[
\boxed{
m_{\text{в-ва, итог}}
=
m_{\text{в-ва, 1}}
+
m_{\text{в-ва, 2}}
}.
\]

Если смешиваются два раствора:

\[
\boxed{
(m_1+m_2)c
=
m_1c_1+m_2c_2
}.
\]

\keyidea{Количество раствора складывается, количество растворённого вещества также складывается. Конечная концентрация определяется их отношением.}

% ============================================================

\topic{9. Пример: смешивание двух растворов}

Смешали $4$ л $15\%$-го раствора и $6$ л $25\%$-го раствора.

Пусть итоговая концентрация равна $x$.

\begin{center}
\renewcommand{\arraystretch}{1.35}
\begin{tabularx}{0.84\linewidth}{@{} X c c c @{}}
\toprule
& \textbf{Количество раствора} & \textbf{Концентрация} & \textbf{Вещество} \\
\midrule
Первый & $4$ & $0{,}15$ & $4\cdot0{,}15$ \\
Второй & $6$ & $0{,}25$ & $6\cdot0{,}25$ \\
Итог & $10$ & $x$ & $10x$ \\
\bottomrule
\end{tabularx}
\end{center}

Уравнение баланса:

\[
10x=4\cdot0{,}15+6\cdot0{,}25.
\]

\[
10x=0{,}6+1{,}5=2{,}1,
\]

\[
x=0{,}21.
\]

Следовательно,

\[
\boxed{21\%}.
\]

Проверка здравого смысла:

\[
15\%<21\%<25\%.
\]

Итоговая концентрация действительно расположена между исходными концентрациями.

% ============================================================

\topic{10. Одинаковые количества растворов}

Смешали одинаковые количества $13\%$-го и $23\%$-го растворов.

Пусть количество каждого раствора равно $m$, итоговая концентрация --- $x$.

\[
2mx=0{,}13m+0{,}23m.
\]

Поскольку

\[
m>0,
\]

можно сократить обе части уравнения на $m$:

\[
2x=0{,}36,
\]

\[
\boxed{x=0{,}18=18\%}.
\]

Для одинаковых количеств двух растворов можно воспользоваться средней концентрацией:

\[
\boxed{
c_{\text{итог}}
=
\frac{c_1+c_2}{2}
}.
\]

В данном случае:

\[
\frac{13\%+23\%}{2}=18\%.
\]

% ============================================================

\topic{11. Добавление воды}

В чистой воде концентрация растворённого вещества равна

\[
\boxed{0}.
\]

Пусть к $8$ л $15\%$-го раствора добавили $4$ л воды.

Количество растворённого вещества сохраняется:

\[
8\cdot0{,}15=1{,}2.
\]

Общее количество раствора:

\[
8+4=12.
\]

Пусть итоговая концентрация равна $x$:

\[
12x=1{,}2.
\]

\[
x=0{,}1.
\]

Следовательно,

\[
\boxed{10\%}.
\]

\keyidea{При добавлении воды количество растворённого вещества сохраняется, а общее количество раствора увеличивается. Поэтому концентрация уменьшается.}

% ============================================================

\topic{12. Если неизвестны количества двух растворов}

Иногда вводятся две неизвестные величины.

Например, если смешали $x$ кг первого раствора и $y$ кг второго, причём итоговая масса известна:

\[
\boxed{x+y=M}.
\]

Второе уравнение получается из баланса растворённого вещества:

\[
\boxed{
c_1x+c_2y=c_{\text{итог}}M
}.
\]

Получается система:

\[
\begin{cases}
x+y=M,\\
c_1x+c_2y=c_{\text{итог}}M.
\end{cases}
\]

Такой подход особенно удобен, когда требуется найти количества обоих исходных растворов.

% ============================================================

\topic{13. Дополнительные переменные}

Дополнительная буква допустима, если точное значение величины неизвестно.

Например, одинаковые массы можно обозначить так:

\[
m,\qquad m.
\]

Тогда после составления уравнения множитель $m$ часто сокращается.

Перед сокращением проверяется условие

\[
m\ne0.
\]

Для массы раствора по смыслу задачи:

\[
m>0.
\]

\keyidea{Удобная дополнительная переменная упрощает математическую модель и часто исчезает после преобразований.}

% ============================================================

\topic{14. Типичные ошибки}

\begin{enumerate}
    \item \textbf{Смешение времени и производительности.}

    Формулировка «6 деталей за час» задаёт производительность:

    \[
    p=6.
    \]

    \item \textbf{Разные единицы времени.}

    Перед составлением уравнения часы и минуты приводятся к одной единице.

    \item \textbf{Неверный порядок разности времён.}

    Используйте:

    \[
    t_{\text{большее}}-t_{\text{меньшее}}>0.
    \]

    \item \textbf{Сложение времён при совместной работе.}

    Для совместно работающих исполнителей основная связь строится через производительности:

    \[
    p_{\text{общ}}=p_1+p_2.
    \]

    \item \textbf{Проценты оставлены в виде целых чисел.}

    В формулах:

    \[
    15\%=0{,}15.
    \]

    \item \textbf{Концентрации складываются напрямую.}

    Основой служит баланс количества растворённого вещества:

    \[
    m_1c_1+m_2c_2=(m_1+m_2)c.
    \]

    \item \textbf{У воды записана ненулевая концентрация вещества.}

    Для чистой воды:

    \[
    c=0.
    \]

    \item \textbf{Игнорируются ограничения на знаменатели.}

    Если встречается

    \[
    \frac{A}{x},
    \]

    требуется

    \[
    x\ne0.
    \]

    Для времени и производительности обычно дополнительно:

    \[
    x>0.
    \]

    \item \textbf{При квадратном уравнении принимается отрицательный корень.}

    В задачах о времени, массе и производительности отрицательное значение противоречит смыслу величины.
\end{enumerate}

% ============================================================

\topic{15. Быстрый алгоритм на экзамене}

\begin{enumerate}
    \item Определите тип задачи.
    \item Выпишите три основные величины.
    \item Сделайте строки для всех объектов или состояний.
    \item Поставьте $x$ в наиболее удобную клетку.
    \item Заполните формульный столбец.
    \item Проверьте единицы измерения.
    \item Найдите словесную связь:
    \[
    \text{вместе},\quad
    \text{одинаковое время},\quad
    \text{раньше на},\quad
    \text{больше на},\quad
    \text{смешали}.
    \]
    \item Запишите соответствующее уравнение.
    \item Укажите ограничения.
    \item После решения проверьте смысл ответа.
\end{enumerate}

\vspace{0.4em}

\begin{center}
\renewcommand{\arraystretch}{1.4}
\begin{tabularx}{0.95\linewidth}{@{} X X @{}}
\toprule
\textbf{Производительность} & \textbf{Растворы} \\
\midrule
$A=pt$ &
$m_{\text{в-ва}}=mc$ \\

$p=\dfrac{A}{t}$ &
$m_{\text{итог}}=m_1+m_2$ \\

$t=\dfrac{A}{p}$ &
$m_1c_1+m_2c_2=(m_1+m_2)c$ \\

$p_{\text{общ}}=p_1+p_2$ &
$c_{\text{воды}}=0$ \\
\bottomrule
\end{tabularx}
\end{center}

% ============================================================
% ТРЕНИРОВОЧНЫЙ БЛОК
% ============================================================

\newpage
\thispagestyle{fancy}

\begin{center}
{\LARGE\bfseries\color{darkaccent}Тренировочный блок}

\vspace{0.25cm}

{\color{accent!55}\rule{0.76\linewidth}{0.7pt}}

\vspace{0.35cm}

{\small\color{muted}
Составьте математическую модель и получите числовой ответ.}
\end{center}

\subtopic{Средний уровень}

\begin{enumerate}

\item
Один мастер выполняет заказ за $4$ часа, другой --- за $6$ часов.
За какое время они выполнят этот заказ, работая вместе?

\item
Петя отвечает на $8$ вопросов теста в час, Ваня --- на $10$ вопросов в час.
Они одновременно начали выполнять одинаковые тесты.
Петя закончил на $15$ минут позже Вани.
Сколько вопросов было в каждом тесте?

\item
Смешали $5$ л $12\%$-го раствора вещества и $3$ л $28\%$-го раствора того же вещества.
Определите концентрацию получившегося раствора.

\end{enumerate}

\subtopic{Уровень выше среднего}

\begin{enumerate}[resume]

\item
Первая труба наполнила бак объёмом $720$ л, а вторая за такое же время --- бак объёмом $840$ л.
Вторая труба пропускает на $5$ л воды в минуту больше первой.
Найдите производительность первой трубы.

\item
Рабочий должен изготовить $120$ деталей.
Фактически он изготавливал на $4$ детали в час больше запланированного и закончил работу на $1$ час раньше.
Сколько деталей в час он планировал изготавливать?

\item
К $8$ л $15\%$-го раствора добавили $4$ л воды.
Определите концентрацию получившегося раствора.

\item
Смешали одинаковые количества $13\%$-го и $23\%$-го растворов одного вещества.
Определите концентрацию получившейся смеси.

\item
Смешали $30\%$-й и $50\%$-й растворы одного вещества.
Получили $20$ кг $42\%$-го раствора.
Сколько килограммов каждого исходного раствора использовали?

\item
К $20$ л $25\%$-го раствора добавили воду.
После этого концентрация стала равна $10\%$.
Сколько литров воды добавили?

\end{enumerate}

\subtopic{Сложный уровень}

\begin{enumerate}[resume]

\item
Две машины, работая вместе, выполняют заказ за $4$ часа.
Первая машина, работая одна, выполняет тот же заказ на $6$ часов быстрее второй.
За сколько часов каждая машина выполняет заказ отдельно?

\end{enumerate}

% ============================================================
% ОТВЕТЫ
% ============================================================

\newpage
\thispagestyle{fancy}

\begin{center}
{\LARGE\bfseries\color{darkaccent}Ответы}

\vspace{0.25cm}

{\color{accent!55}\rule{0.55\linewidth}{0.7pt}}
\end{center}

\vspace{0.8cm}

\begin{table}[ht]
\centering
\caption{Ответы к тренировочному блоку}
\renewcommand{\arraystretch}{1.45}
\begin{tabularx}{0.82\linewidth}{@{} c X @{}}
\toprule
\textbf{№} & \textbf{Ответ} \\
\midrule

1 &
$2{,}4$ ч, то есть $2$ ч $24$ мин \\

2 &
$10$ вопросов \\

3 &
$18\%$ \\

4 &
$30$ л/мин \\

5 &
$20$ деталей в час \\

6 &
$10\%$ \\

7 &
$18\%$ \\

8 &
$8$ кг $30\%$-го раствора и $12$ кг $50\%$-го раствора \\

9 &
$30$ л воды \\

10 &
первая машина --- $6$ ч, вторая --- $12$ ч \\

\bottomrule
\end{tabularx}
\end{table}

\vfill

\begin{center}
{\small\color{muted}
Перед сдачей ответа проверяйте единицы измерения,
положительность времени, массы и производительности,
а также перевод концентрации из долей в проценты.}
\end{center}

\end{document}